% Generated from locale/tr/content/intuitionistic-logic/introduction/axiomatic-derivations.tex; do not hand-edit.


\olfileid[tr]{int}{int}{axd}

\olsection{Aksiyomatik \usetoken{P}{derivation}}

Sezgici önerme mantığı için aksiyomatik !!{derivation}, kavramsal bakımdan en
basit ve tarihsel olarak ilk !!{derivation} sistemleridir. Klasik önerme
mantığındakiyle aynı biçimde işlerler.

\begin{defn}[!!^{derivability}]
$\Gamma$, $\Lang L$ dilinin !!{formula}s{}den oluşan bir kümeyse bir
\emph{!!{derivation}}, $\Gamma$'dan elde edilen sonlu bir $!A_1$, \dots,~$!A_n$ !!{formula}
dizisidir ve her $i \le n$ için aşağıdakilerden biri geçerlidir:
\begin{enumerate}
\item $!A_i \in \Gamma$; ya da
\item $!A_i$ bir aksiyomdur; ya da
\item $!A_i$, bazı $!A_j$ ve $!A_k$'dan modus ponens yoluyla çıkar; burada
  $j < i$ ve $k < i$'dir, yani $!A_k \ident !A_j \lif !A_i$.
\end{enumerate}
\end{defn}

\begin{defn}[Aksiyomlar]
Sezgici önerme mantığının $\PAx$ \emph{aksiyomlar} kümesi, aşağıdaki
biçimlerdeki bütün !!{formula}s{}i içerir:
\begin{align}
  & (!A \land !B) \lif !A \ollabel{ax:land1}\\
  & (!A \land !B) \lif !B \ollabel{ax:land2}\\
  & !A \lif (!B \lif (!A \land !B)) \ollabel{ax:land3}\\
  & !A \lif (!A \lor !B) \ollabel{ax:lor1}\\
  & !A \lif (!B \lor !A) \ollabel{ax:lor2}\\
  & (!A \lif !C) \lif ((!B \lif !C) \lif ((!A \lor !B) \lif !C)) \ollabel{ax:lor3}\\
  & !A \lif (!B \lif !A) \ollabel{ax:lif1}\\
  & (!A \lif (!B \lif !C)) \lif ((!A \lif !B) \lif (!A \lif !C)) \ollabel{ax:lif2}\\
  & \lfalse \lif !A \ollabel{ax:lfalse1}
\end{align}
\end{defn}

\begin{defn}[!!^{derivability}]
!!^a{formula}~$!A$'nın \emph{!!{derivable}} olması $\Gamma$'ya göre
tanımlanır. Bu bağıntıyı $\Gamma \Proves !A$ diye yazarız; bağıntı,
$\Gamma$'dan $!A$ ile biten !!a{derivation} varsa geçerlidir.
\end{defn}

\begin{defn}[Teoremler]
!!^a{formula}~$!A$, boş kümeden~$!A$'nın !!a{derivation}{}i varsa bir
\emph{teorem}dir. $\Proves !A$ yazımı, $!A$ bir teoremse kullanılır; değilse
$\Proves/ !A$ yazarız.
\end{defn}

\begin{prop}
  Sezgici mantığın aksiyomatik !!{derivation} sisteminde $\Gamma \Proves !A$
  ise klasik mantıkta da $\Gamma \Proves !A$ olur. Özellikle $!A$ sezgici bir
  teoremse klasik bir teoremdir.
\end{prop}

\begin{proof}
  Her sezgici aksiyom aynı zamanda klasik bir aksiyomdur; dolayısıyla sezgici
  mantıktaki her !!{derivation}, klasik mantıkta da !!a{derivation}{}dir.
\end{proof}

